\scanpage{166}
Pour définir un hom $F$, il faut \unsure{définir}
\diagram{\begin{tikzcd}[ampersand replacement=\&,row sep=large,column sep=small]
\& (E_1\otimes L_2)\unsure{\oplus|\amalg}(L_1\otimes E_2)=\Sigma
\arrow[dl,"{\mathrm{id}_1\otimes\alpha_2,\ \alpha_1\otimes\mathrm{id}_2}"']
\arrow[dr,"{\pi_1\otimes\mathrm{id}_{\unsure{2|L_2}},\ \mathrm{id}_1\otimes\pi_2}"]\&\\
E_1\otimes E_2\arrow[d,"F_{11}"']\&\&L_1\otimes L_2\arrow[d,"F_{22}"]\arrow[dll,"F_{21}" near start]\\
E_1\otimes\deleted{\unsure{E_2\otimes}}E_2
\arrow[dr,"{\mathrm{id}_1\otimes\psi_2,\ \psi_1\otimes\mathrm{id}_2}"']\&\&
M_1\otimes M_2\arrow[dl,"{\unsure{\varpi|\overline{\omega}}_1\otimes\mathrm{id}_2,\ \mathrm{id}_1\otimes\varpi_2}"]\\
\&(E_1\otimes M_2)\times(M_1\otimes E_2)=\Pi\arrow[from=2-1,to=3-3,crossing over,"F_{12}" near end]\&
\end{tikzcd}}
\unsure{avec|car} 4 carrés commutatifs

Si on veut que $F$ induise \deleted{l'identité}\inserted{un hom d'extensions, donc} $L_1\otimes L_2\longrightarrow L_1\otimes L_2$, il f. et s. \deleted{\unsure{avoir}} que $\boxed{F_{22}=0}$, si on veut \inserted{de plus} que \unsure{ça|$\Sigma$} induise l'identité sur $L_1\otimes L_2$, il f. et s. que $\boxed{F_{21}=\alpha_1\otimes\alpha_2}$, \unsure{d\textsuperscript{o}|si|et} \unsure{si|{}} on veut que \unsure{ça|$F$} induise l'identité sur $M_1\otimes M_2$, il f. et s. que $\boxed{F_{12}=\psi_1\otimes\psi_2}$ : \deleted{\unsure{Si|D}} ces conditions sont vérifiées, alors \deleted{\unsure{d}}\unsure{deux} des 4 carrés voulus sont déjà commutatifs, savoir
\diagram{\begin{tikzcd}[ampersand replacement=\&,column sep=small,row sep=small]
\&\Sigma\arrow[dl,dash]\arrow[dr,dash]\&\\
E_1\otimes E_2\arrow[rrd,"0","{F_{12}=\psi_1\otimes\psi_2}"']\&\&L_1\otimes L_2\arrow[d,"{F_{22}=0}"]\\
\&\&M_1\otimes M_2
\end{tikzcd}}
et
\diagram{\begin{tikzcd}[ampersand replacement=\&,column sep=small,row sep=small]
\&\&L_1\otimes L_2\arrow[dll,"{\deleted{E_1\otimes E_2}F_{21}}"']\arrow[d,"{F_{22}=0}"]\\
E_1\otimes E_2\arrow[dr]\&\&M_1\otimes M_2\arrow[dl]\\
\&\deleted{\unsure{\varphi}}\Pi\&
\end{tikzcd}}
\margin{(les quatre composés sont nuls)}
% archive sheet number: 165
